\documentclass[11pt]{article}
\usepackage[margin=1in]{geometry}
\usepackage{amsmath,amssymb,booktabs}
\usepackage[hidelinks]{hyperref}
\newcommand{\HH}{\mathrm{HH}}
\newcommand{\HV}{\mathrm{HV}}
\newcommand{\kma}{\mathrm{km}^{2}}
\title{Dancing with the SARs\\Radar Detection and Validation}
\author{Science review for Abel and Harutyun}
\date{}
\begin{document}
\maketitle

\section{What the figures show}

The Tonle Sap change-date map shows earlier changes in parts of the
southeast corridor and later changes around parts of the lake margin and
northwest. This is consistent with seasonal reverse flow from the Mekong.
The map colors give the first observation of a new radar brightness level.
Rain, crops, mixed pixels, and swath-edge effects can also change brightness.

The open-water class expands. The area flagged as possible flooded
vegetation peaks on August 30 and then declines. Both radar channels
darkened at Prek Toal. Reduced vegetation contrast as water expands is a
possible explanation. Use \emph{Possible flooded vegetation} in the legend.

\section{Radar brightness and valid pixels}

Smooth water often reflects radar away from the satellite. Wind can roughen
water and increase the return. Water and upright trunks can produce a double
reflection. HH transmits and receives horizontal polarization. HV transmits
horizontal polarization and receives vertical polarization.

For aligned preview brightness $B_t(p)\in[0,1]$ at pixel $p$ and date $t$,
\begin{equation}
 \Delta B(p)=B_{\mathrm{after}}(p)-B_{\mathrm{before}}(p).
\end{equation}
A change from $0.06$ to $0.12$ is $+0.06$ display units. Preview brightness
uses an unknown display stretch. With positive calibrated linear power $P$,
the physical backscatter change is
\begin{equation}
 \Delta_{\mathrm{dB}}=10\log_{10}\left(
 \frac{P_{\mathrm{after}}}{P_{\mathrm{before}}}\right).
\end{equation}
Doubling calibrated power gives approximately $+3.01\,\mathrm{dB}$.
Average calibrated power in linear units before converting to decibels.

Let $V(q)=1$ for a valid pixel and $0$ otherwise. The mean over a
neighborhood $W(p)$ is
\begin{equation}
 m_t(p)=\frac{\sum_{q\in W(p):\,V(q)=1} B_t(q)}
              {\sum_{q\in W(p)}V(q)}.
\end{equation}
Exclude windows with zero valid pixels. The $3\times3$ filter spans roughly
$540\,\mathrm{m}$, so shorelines and small fields can mix surface types.

\section{Open water and possible flooded vegetation}

For processed preview brightness in the two channels, define seeds and
candidates on the common valid footprint:
\begin{align}
 S_t &= \{p:V(p)=1,\ B_{\HH,t}(p)<0.06,\ B_{\HV,t}(p)<0.35\},\\
 C_t &= \{p:V(p)=1,\ B_{\HH,t}(p)<0.12,\ B_{\HV,t}(p)<0.35\}.
\end{align}
Keep candidates connected to a seed by an eight-neighbor path within $C_t$,
then remove tiny connected components. HH $=0.09$ and HV $=0.20$ qualifies
if the pixel connects to a seed and survives cleanup. A connected dark
field can also pass.

For channel $c$, calculate change from the mean of the first two dates:
\begin{equation}
 \Delta B_{c,t}(p)=B_{c,t}(p)-\frac{B_{c,1}(p)+B_{c,2}(p)}{2}.
\end{equation}
The vegetation rule requires $\Delta B_{\HH,t}>0.15$ and
$\Delta B_{\HV,t}<0.06$, excluding current open water and likely permanent
water. The peak flagged area is sensitive to the HH cutoff.
\begin{center}
\begin{tabular}{rr}
\toprule
HH increase cutoff & Peak flagged area ($\kma$)\\
\midrule
0.12 & 1,083\\
0.15 & 611\\
0.20 & 139\\
\bottomrule
\end{tabular}
\end{center}

To check whether the same pixels later look like open water, let
$F_{\mathrm{Aug30}}$ be the August vegetation mask and
$M_{\mathrm{Sep23}}$ the September open-water mask. On the common
footprint, for a nonempty August mask, calculate
\begin{equation}
 f=\frac{|F_{\mathrm{Aug30}}\cap M_{\mathrm{Sep23}}|}
         {|F_{\mathrm{Aug30}}|}.
\end{equation}
Compare $f$ across thresholds and against gauge and optical observations.
This fraction is a proposed check with the result still unknown.

\section{First observed change}

For a candidate split with $n_1$ earlier and $n_2$ later observations,
\begin{equation}
 \Delta=\overline B_{\mathrm{after}}-\overline B_{\mathrm{before}},
 \qquad S=|\Delta|\sqrt{\frac{n_1n_2}{n_1+n_2}}.
\end{equation}
Five observations on either side and a $0.20$-unit step give
$S\approx0.316$. The detector screens for minimum step size, persistence,
and a $t$-like statistic. Temporal dependence and the search across
candidate splits limit statistical interpretation. Report the first
observation of the new state and the interval since the preceding pass.

\section{Area and validation}

For a classified mask $M$ with pixel ground areas $a_p$ on the common
footprint,
\begin{equation}
 A(M)=\sum_{p\in M}a_p,
 \qquad a_p\approx(0.18\,\mathrm{km})^2=0.0324\,\kma
\end{equation}
for nominal square $180\,\mathrm{m}$ cells. A uniform shoreline offset
$d$ along perimeter $L$ gives $|\Delta A|\approx L|d|$. A one-pixel offset
around a $250\,\mathrm{km}$ shoreline would affect about $45\,\kma$.

For a bank at angle $\theta$ above horizontal,
$|\Delta x|\approx|\Delta h|/\tan\theta$. A $0.12\,\mathrm{m}$ drop moves
the waterline about $6.9\,\mathrm{m}$ at $1^\circ$ or $69\,\mathrm{m}$
at $0.1^\circ$. These are illustrative slopes.

For classified water $M$ and reference water $R$ on matching grids,
\begin{equation}
 \mathrm{IoU}=\frac{|M\cap R|}{|M\cup R|},
 \qquad \text{share inside reference}=\frac{|M\cap R|}{|M|}.
\end{equation}
The team reports that $98$--$99\%$ of Tonle Sap late-season classified
water lies inside the historical maximum extent, and that Sevan permanent
water has IoU $0.94$. About $97\%$ of Ararat's latest classified water lies
outside the historical reference. The
\href{https://global-surface-water.appspot.com/}{JRC Landsat reference}
ends in 2021, so this checks spatial agreement. The team's deep-learning
perturbation tests measure stability under simulated noise relative to
rule-generated masks or clean-input predictions.

\section{Numbers for the demo}

Use the same detector output for the story, map, and chart. The math run
classifies about $2{,}234\,\kma$ on June 19 and $3{,}837\,\kma$ on
October 17, an increase of roughly $72\%$ within the common footprint.
The exploratory brightness analysis gives different extents.

\begin{center}
\begin{tabular}{lrrr}
\toprule
2026 observation interval & Change ($\kma$) & Days & $\kma$/day\\
\midrule
July 25 to August 18 & 381 & 24 & 16\\
September 11 to September 23 & 375 & 12 & 31\\
\bottomrule
\end{tabular}
\end{center}
Each rate is the net classification difference divided by elapsed days.
Within-interval rates are unknown. For the story, use \emph{marked
expansion between September 11 and 23}.

\paragraph{Before publishing.}
Verify the sampling circles against higher-resolution imagery and valid
radar coverage. Recompute the charts for the five moved centers. Show the
Tonle Sap observation gap, use Ararat to explain the preview limits, and
check the deployed figures, public links, and AI-use statement.
\end{document}
